% Plan de tesis -- Capitulo 2.
%
% Correlato en el plan de ejemplo: capitulo 2 (Materiales y Metodos), con una seccion de base
% tecnica sobre la que se apoya la propuesta, una de Metodos (que es donde el ejemplo enuncia su
% hipotesis) y una de Datos. Extension del ejemplo: ~4 paginas.
\chapter{Materiales y Métodos}

\section{Funciones de pérdida asimétricas}

% Base tecnica sobre la que se apoya la propuesta. Las dos familias y sus referencias canonicas ya
% estan en el catalogo (koenker-bassett-1978, newey-powell-1987, ehm-2016-quantiles-expectiles).
Sea $u = y - \hat{y}$ el error de pronóstico y $\tau \in (0,1)$ un parámetro de asimetría.
La regresión cuantílica \cite{koenker-bassett-1978} minimiza la pérdida \textit{pinball}
\begin{equation}
  \rho^{\mathrm{q}}_{\tau}(u) = u \left( \tau - \mathbf{1}\{u < 0\} \right),
\end{equation}
cuyo minimizador poblacional es el cuantil de nivel $\tau$ de la distribución de $y$.
La regresión expectílica \cite{newey-powell-1987} reemplaza el valor absoluto implícito por un
término cuadrático,
\begin{equation}
  \rho^{\mathrm{e}}_{\tau}(u) = \left| \tau - \mathbf{1}\{u < 0\} \right| u^{2},
\end{equation}
y su minimizador es el expectil de nivel $\tau$.
Para $\tau = 1/2$ ambas se reducen a los casos simétricos conocidos: el error absoluto y el error
cuadrático.
Con $\tau > 1/2$ la pérdida penaliza más la subestimación que la sobreestimación.

La diferencia práctica entre las dos familias es la ponderación de los errores grandes.
El expectil conserva el término cuadrático, de modo que el gradiente crece con la magnitud del
error; eso lo hace derivable en todo el dominio y directamente utilizable como objetivo de
entrenamiento por descenso de gradiente, a diferencia de la pérdida \textit{pinball}, que no es
derivable en el origen.

% PENDIENTE: cerrar esta seccion con la discusion de elicitabilidad --- que funcional estima cada
% perdida y por que eso condiciona la comparacion entre modelos --- cuando este consolidado el
% relevamiento (eje b).

\section{Métodos}

% En el plan de ejemplo esta es la seccion donde se enuncia la hipotesis del trabajo. Se replica.

% PROVISORIO: la hipotesis la cierra el autor con el director.
La hipótesis del trabajo es que el nivel de asimetría conveniente no es constante: en estiaje el
costo de sobreestimar y el de subestimar son comparables, mientras que en crecida la asimetría es
marcada.
Se propone entonces modular el parámetro de asimetría con el régimen hidrológico, es decir,
reemplazar el $\tau$ fijo por
\begin{equation}
  \tau_t = g(\mathbf{x}_t),
\end{equation}
donde $\mathbf{x}_t$ resume el estado de la cuenca en el instante $t$ y $g$ es un modulador
acotado en $(0,1)$.

% PENDIENTE: la forma funcional concreta de g, el conjunto de modelos a comparar y el protocolo de
% validacion se completan con lo que devuelva el relevamiento. Lo ya implementado en el repositorio
% (validacion walk-forward, linea de base de persistencia, contraste de Diebold-Mariano) se
% describe en el capitulo 3.

\section{Datos}

% Inventario verificado contra la documentacion del repositorio (docs/roadmap.md,
% docs/data_sources.md, docs/gold_consolidation_contract.md, docs/gold_quality_report.md).
% Las cifras corresponden a la medicion del 2026-08-26 y hay que refrescarlas antes de elevar.
El conjunto de datos ya está construido y es el insumo de esta tesis.
Cubre la cuenca alta del río Uruguay, aguas arriba de la frontera entre Brasil y Argentina, y se
consolida con una arquitectura por capas sobre \textit{Databricks}.

La variable objetivo es el caudal diario, en m\textsuperscript{3}/s, en la estación 74100000 (Irai)
de la Agência Nacional de Águas de Brasil, sobre la frontera.
El caudal no se mide directamente: se deriva del nivel observado mediante la curva de aforo de la
estación, una ley de potencia $Q = A\,(H - H_0)^{N}$ cuyos coeficientes cambian por tramo y por
período de vigencia.
El nivel se conserva como objetivo secundario, de modo que la serie medida nunca se pierde detrás
de la conversión.
Los horizontes de pronóstico son t+1 a t+7, más t+14.

\begin{table}[H]
  \centering
  \caption{Fuentes de datos del conjunto consolidado.}
  \label{tab:fuentes}
  \small
  \begin{tabular}{@{}llll@{}}
    \toprule
    \textbf{Fuente} & \textbf{Variable} & \textbf{Tipo} & \textbf{Cobertura} \\
    \midrule
    ANA (Brasil)        & Nivel y caudal      & Observado   & desde 1950 \\
    ANA (Brasil)        & Precipitación       & Observado   & desde 1923 \\
    ANA (Brasil)        & Curvas de aforo     & Observado   & desde 1948 \\
    CPTEC/INPE MERGE    & Precipitación       & Observado, grilla 0,1\textdegree{} & desde 1998 \\
    CPTEC/INPE SAMeT    & Temperatura         & Observado, grilla 0,05\textdegree{} & desde 2000 \\
    INMET (Brasil)      & Temperatura         & Observado   & desde 2000 \\
    GEFS Reforecast v12 & Precipitación       & Pronosticado & 2000--2019 \\
    ECMWF TIGGE         & Precipitación       & Pronosticado & desde 2006 \\
    \bottomrule
  \end{tabular}
\end{table}

La distinción entre las dos últimas filas y el resto es la que gobierna el diseño del conjunto: las
variables pronosticadas son las únicas que estarán disponibles en el momento en que el modelo deba
operar, y por eso el histórico de pronósticos --- y no el de observaciones --- es el que fija el
piso temporal del conjunto en el año 2000.
El resultado es una tabla de 9.732 días entre el 1 de enero de 2000 y el 23 de agosto de 2026, con
83 columnas, sin huecos de calendario.
El agregado de precipitación y caudal de la cuenca se construye sobre 36 estaciones.

% FALTA: refrescar las cifras contra Databricks antes de elevar el plan, y decidir con el director
% si conviene declarar las limitaciones de calidad ya documentadas (curvas de aforo extrapoladas en
% algunas estaciones del agregado, desfasaje de la ventana diaria de MERGE contra la de ANA).
